\section{解決策としての提案手法}
\label{sec:proposal}

\subsection{提案の概要}
\label{subsec:overview}

前章の主張を実現するため，本提案は次の 3 種類のノードからなる構成を取る．

\begin{itemize}
  \item \textbf{指揮者マイコン}（Arduino UNO R4 WiFi~\cite{arduino_uno_r4}）$\times\,1$ 台
  \item \textbf{楽器マイコン}（Arduino UNO R4 WiFi）$\times\,4$ 台
  \item \textbf{PC}（Processing による音響合成）$\times\,1$ 台
\end{itemize}

指揮者マイコンは内蔵 IMU で指揮棒の動作を読み取り，拍・テンポ・強弱などの演奏パラメータを推定して，WiFi 経由の UDP ブロードキャストで全楽器マイコンへ同時配信する．楽器マイコンは受信した拍情報に従って音符データを PC へ送出し，PC 側は Processing で金管楽器などの音色を合成・出力する．

本稿は筆者が担当する指揮者マイコン部の設計方針に焦点を絞り，IMU によるジェスチャ取得（\ref{subsec:imu_gesture}）と WiFi/UDP ブロードキャストによる配信（\ref{subsec:wifi_udp}）という 2 つの技術選択を中心に述べる．

\subsection{IMU によるジェスチャ取得}
\label{subsec:imu_gesture}

指揮動作のセンシングには，Arduino UNO R4 WiFi に内蔵されている 6 軸 IMU（LSM6DSOX~\cite{lsm6dsox}）を用いる．外付けセンサや追加配線を要さないため，指揮棒型の筐体に収める際の機構負担が小さい．

加速度と角速度を 100\,Hz 以上でサンプリングし，ノイズ除去後に加速度ノルムのピーク検出によって拍タイミングを抽出する．拍間のインターバルを移動平均してテンポを推定し，ピーク時の加速度振幅から強弱パラメータを導出する．本提案ではまず加速度ノルムに基づく簡便なアルゴリズムから着手し，必要に応じて閾値やフィルタの調整，あるいは学習ベース手法への置き換えを検討する．

\subsection{WiFi/UDP ブロードキャストによる配信}
\label{subsec:wifi_udp}

指揮者マイコンから楽器マイコン群への拍情報配信には，WiFi 上の UDP ブロードキャストを用いる．UNO R4 WiFi は ESP32-S3 を無線コプロセッサとして内蔵しているため，追加ハードウェアなしで同一サブネット上の全ノードに対し単一パケットで同時に送信できる．

UDP はコネクションレスで再送を行わないため遅延が小さく，1 対多の同時配信に適する．パケットロスが発生した場合も拍は周期的に送出されるため，次拍で自然に復帰でき，合奏の継続性を損ないにくい．合奏におけるアンサンブル精度の観点からは，片道遅延はできるだけ小さく抑えることが望ましく，本提案では次節で数値目標を設定する．

\subsection{指揮者マイコンの要件}
\label{subsec:requirements}

以上の技術選択を踏まえ，指揮者マイコン部に対して次の要件を設定する．

\begin{itemize}
  \item IMU サンプリング周波数: 100\,Hz 以上（手振りの最短周期に対して十分な分解能）
  \item 指揮動作から拍検出までの処理遅延: 50\,ms 以下
  \item 指揮者マイコンから楽器マイコンまでの UDP 片道遅延: 30\,ms 以下
\end{itemize}

これらは合奏で許容される同期ズレを入力〜検出〜配信の各段階に配分した目安であり，実機での実現可能性は次週以降の個人タスクにおいて計測・評価する．

\subsection{処理方針}
\label{subsec:processing}

指揮者マイコンの処理は，入力・ロジック・出力の 3 段ループで構成する．

\begin{enumerate}
  \item \textbf{入力}: IMU から加速度・角速度を一定周期で取得する
  \item \textbf{ロジック}: ノイズ除去後，加速度ノルムのピーク検出により拍を抽出し，拍間隔からテンポと強弱を推定する
  \item \textbf{出力}: 拍タイミングとテンポ・強弱を UDP ブロードキャストで楽器マイコンへ送出する
\end{enumerate}

この 3 段構成は筆者が過去に採用している Embedded-Module-Architecture~\cite{embedded_module_arch} に準拠する．段ごとのモジュール境界を明確にすることで，拍検出アルゴリズムの差し替えや，入力を IMU からシミュレータに切り替えた単体テストが容易になる．
